\documentclass[11pt, ukrainian]{article}
\setlength\textwidth{180mm} \setlength\textheight{277mm}
\setlength\voffset{-38mm} \setlength\hoffset{-28mm}
%\setlength\parindent{0mm}
\pagestyle{empty}
\usepackage[T2A]{fontenc}
\usepackage[utf8]{inputenc}
\usepackage[ukrainian]{babel}
\begin{document}
{\Large \center  Варіанти завдань до лабораторної роботи №4\par 2020/21 н.р.\par \bigskip\bigskip}

Порядок полів у записах та інформацію додаткового файлу фіксовано у варіанті.

Для порівняння числових значень використовується традиційний порядок.

Для порівняння рядкових значень використовується лексикографічний порядок. 

У випадку сортування за кількома полями результуюче сортування будується 
за порядками окремих полів за допомогою конструкції лексикографічного порядку. 
Послідовність полів, за якими має відбуватися сортування, задається в умові. 
Якщо не вказано інше, то поле сортується за зростанням.

Якщо умова вимагає знайти $k$ най… і $k$-е місце займають кілька, то слід виводити всіх їх.

\newpage
\par \bigskip

{\center Варіант  \No \textbf { 1 }\par}
\bigskip
%type = booking
%variant = 124
Обробляється інформація щодо відвантажених накладних.

\underline{Записи основного файлу містять поля}: номер накладної, кількість, назва товару, вартість, ціна, номер товарної позиції у межах накладної.

\underline{У допоміжному файлі наявні ключі}: найбільша ціна товарної позиції, сума вартостей.

\underline{Знайти} 5 найдорожчих товарів (з найбільшою середньою ціною відвантаження). Вивести по кожному з них інформацію:\par
{\noindent}--- на першому рядку:\par
назва товару, загальна відвантажена кількість, середня ціна відвантаження (з 2 знаками після крапки), найменша ціна відвантаження,  найбільша ціна відвантаження;

{\noindent}--- на наступних рядках, починаючи з табуляції, вивести для товару всі накладні, де він відвантажувався за ненайбільшою своєю ціною (по одному на рядок):\par
ціна, кількість, ціна, вартість, номер накладної

{\noindent}у такому сортуванні: вартість (за спаданням), ціна(за спаданням), номер накладної.

%Товари сортувати: середня ціна відвантаження (за спаданням), назва товару.

\bigskip\textit{Обмеження}

Кількість, ціна та вартість є значеннями дійсного типу.
Решта полів є непорожніми рядками.
Рядки починатися та закінчуватися білими символами не можуть.

Номер накладної складається з 12 символів (цифри, латинські літери).
Назва товару може мати від 1 до 29 символів включно.
Може містити тільки цифри, літери, пробіли, апостроф, дефіс, підкреслення.

Кількість, ціна мають бути додатними.
Кількість вказується з рівно 4 знаками після крапки.
Ціна вказується з рівно 4 знаками після крапки.
Вартість має бути узгоджена з кількістю та ціною та записується рівно з двома десятковими знаками.
(Застосовується округлення.)

Товари однозначно ідентифікуються назвами та одиницями виміру (за наявності).

Одиниця виміру (за наявності у варіанті) може приймати значення: шт., кг, м , кор.

Кожна накладна має свій власний унікальний номер.
У накладній можуть відвантажуватися кілька товарів (не менше 1).
Товарна позиція --- рядок накладної: товар, кількість, ціна, вартість та, за наявності, одиниця виміру.
Товарні позиції кожної накладної нумеруються послідовними натуральними числами, починаючи з 1.
В одній накладній товари можуть повторюватися, 
тобто товарні позиції однієї накладної можуть містити однакові товари 
(при цьому ціна може бути і різною, і однаковою).

\bigskip\textit{Для деяких варіантів може знадобитись}

Недоотриману вигоду розраховувати як суму виразів
{\center (найбільша ціна відвантаження товару – ціна відвантаження товару)*кількість\par}

\bigskip
{\noindent} округлених до 2 знаків після крапки.


\par
\newpage
{\center Варіант  \No \textbf { 2 }\par}
\bigskip
%type = booking
%variant = 125
Обробляється інформація щодо відвантажених накладних.

\underline{Записи основного файлу містять поля}: кількість, одиниця виміру, ціна, назва товару, номер товарної позиції у межах накладної, вартість, номер накладної.

\underline{У допоміжному файлі наявні ключі}: кількість записів, кількість різних товарів.

\underline{Знайти} 7 найдешевших товарів (з найменшою середньою ціною відвантаження). Вивести по кожному з них інформацію:\par
{\noindent}--- на першому рядку:\par
назва товару, одиниця виміру, кількість накладних, де відвантажувався товар, кількість накладних, де він відвантажувався за своєю максимальною ціною, середня ціна відвантаження (з 2 знаками після крапки);

{\noindent}--- на наступних рядках, починаючи з табуляції, вивести для товару всі накладні, де він відвантажувався за своєю мінімальною ціною (по одному на рядок):\par
номер накладної, ціна, кількість, ціна, вартість

{\noindent}у такому сортуванні: номер накладної, ціна.

%Товари сортувати: кількість накладних, де відвантажувався товар, назва товару, одиниця виміру.

\bigskip\textit{Обмеження}

Кількість, ціна та вартість є значеннями дійсного типу.
Решта полів є непорожніми рядками.
Рядки починатися та закінчуватися білими символами не можуть.

Номер накладної складається з 14 символів (цифри, латинські літери).
Назва товару може мати від 4 до 25 символів включно.
Може містити тільки цифри, літери, пробіли, апостроф, дефіс, підкреслення.

Кількість, ціна мають бути додатними.
Кількість вказується з рівно 3 знаками після крапки.
Ціна вказується з рівно 3 знаками після крапки.
Вартість має бути узгоджена з кількістю та ціною та записується рівно з двома десятковими знаками.
(Застосовується округлення.)

Товари однозначно ідентифікуються назвами та одиницями виміру (за наявності).

Одиниця виміру (за наявності у варіанті) може приймати значення: шт., кг, м , кор.

Кожна накладна має свій власний унікальний номер.
У накладній можуть відвантажуватися кілька товарів (не менше 1).
Товарна позиція --- рядок накладної: товар, кількість, ціна, вартість та, за наявності, одиниця виміру.
Товарні позиції кожної накладної нумеруються послідовними натуральними числами, починаючи з 1.
В одній накладній товари можуть повторюватися, 
тобто товарні позиції однієї накладної можуть містити однакові товари 
(при цьому ціна може бути і різною, і однаковою).

\bigskip\textit{Для деяких варіантів може знадобитись}

Недоотриману вигоду розраховувати як суму виразів
{\center (найбільша ціна відвантаження товару – ціна відвантаження товару)*кількість\par}

\bigskip
{\noindent} округлених до 2 знаків після крапки.


\par
\newpage
{\center Варіант  \No \textbf { 3 }\par}
\bigskip
%type = booking
%variant = 117
Обробляється інформація щодо відвантажених накладних.

\underline{Записи основного файлу містять поля}: ціна, номер накладної, назву товару, вартість, одиниця виміру, номер товарної позиції у межах накладної, кількість.

\underline{У допоміжному файлі наявні ключі}: сума цін, найменша ціна товарної позиції.

\underline{Знайти} 6 найменш популярних товарів (ті, що відвантажувалися в найменшій кількості накладних). Вивести по кожному з них інформацію:\par
{\noindent}--- на першому рядку:\par
назва товару, одиниця виміру, загальна відвантажена кількість, найменша ціна відвантаження (з 2 знаками після крапки);

{\noindent}--- на наступних рядках, починаючи з табуляції, вивести для товару всі накладні, де він відвантажувався за найбільшою ціною (по одному на рядок):\par
ціна, кількість, номер накладної

{\noindent}у такому сортуванні: ціна, номер накладної.

%Товари сортувати: загальна відвантажена кількість, найменша ціна відвантаження, назва товару, одиниця виміру.

\bigskip\textit{Обмеження}

Кількість, ціна та вартість є значеннями дійсного типу.
Решта полів є непорожніми рядками.
Рядки починатися та закінчуватися білими символами не можуть.

Номер накладної складається з 15 символів (цифри, латинські літери).
Назва товару може мати від 3 до 24 символів включно.
Може містити тільки цифри, літери, пробіли, апостроф, дефіс, підкреслення.

Кількість, ціна мають бути додатними.
Кількість вказується з рівно 2 знаками після крапки.
Ціна вказується з рівно 2 знаками після крапки.
Вартість має бути узгоджена з кількістю та ціною та записується рівно з двома десятковими знаками.
(Застосовується округлення.)

Товари однозначно ідентифікуються назвами та одиницями виміру (за наявності).

Одиниця виміру (за наявності у варіанті) може приймати значення: шт., кг, м , кор.

Кожна накладна має свій власний унікальний номер.
У накладній можуть відвантажуватися кілька товарів (не менше 1).
Товарна позиція --- рядок накладної: товар, кількість, ціна, вартість та, за наявності, одиниця виміру.
Товарні позиції кожної накладної нумеруються послідовними натуральними числами, починаючи з 1.
В одній накладній товари можуть повторюватися, 
тобто товарні позиції однієї накладної можуть містити однакові товари 
(при цьому ціна може бути і різною, і однаковою).

\bigskip\textit{Для деяких варіантів може знадобитись}

Недоотриману вигоду розраховувати як суму виразів
{\center (найбільша ціна відвантаження товару – ціна відвантаження товару)*кількість\par}

\bigskip
{\noindent} округлених до 2 знаків після крапки.


\par
\newpage
{\center Варіант  \No \textbf { 4 }\par}
\bigskip
%type = booking
%variant = 122
Обробляється інформація щодо відвантажених накладних.

\underline{Записи основного файлу містять поля}: назва товару, кількість, ціна, номер товарної позиції у межах накладної, номер накладної, вартість, одиниця виміру.

\underline{У допоміжному файлі наявні ключі}: кількість накладних, кількість записів.

\underline{Знайти} 5 найдорожчих товарів (з найбільшою ціною відвантаження). Вивести по кожному з них інформацію:\par
{\noindent}--- на першому рядку:\par
найбільша ціна відвантаження, найменша ціна відвантаження, назва товару, одиниця виміру, загальна відвантажена кількість;

{\noindent}--- на наступних рядках, починаючи з табуляції, вивести для товару всі накладні, де він відвантажувався за ціною, що відрізняється від найбільшої (по одному на рядок):\par
номер накладної, кількість, ціна, недоотримана вигода (з 1 знаком після крапки)

{\noindent}у такому сортуванні: номер накладної.

%Товари сортувати: різниця максимальної та мінімальної відпускних цін за спаданням, назва товару, одиниця виміру.

\bigskip\textit{Обмеження}

Кількість, ціна та вартість є значеннями дійсного типу.
Решта полів є непорожніми рядками.
Рядки починатися та закінчуватися білими символами не можуть.

Номер накладної складається з 11 символів (цифри, латинські літери).
Назва товару може мати від 2 до 28 символів включно.
Може містити тільки цифри, літери, пробіли, апостроф, дефіс, підкреслення.

Кількість, ціна мають бути додатними.
Кількість вказується з рівно 1 знаками після крапки.
Ціна вказується з рівно 4 знаками після крапки.
Вартість має бути узгоджена з кількістю та ціною та записується рівно з двома десятковими знаками.
(Застосовується округлення.)

Товари однозначно ідентифікуються назвами та одиницями виміру (за наявності).

Одиниця виміру (за наявності у варіанті) може приймати значення: шт., кг, м , кор.

Кожна накладна має свій власний унікальний номер.
У накладній можуть відвантажуватися кілька товарів (не менше 1).
Товарна позиція --- рядок накладної: товар, кількість, ціна, вартість та, за наявності, одиниця виміру.
Товарні позиції кожної накладної нумеруються послідовними натуральними числами, починаючи з 1.
В одній накладній товари можуть повторюватися, 
тобто товарні позиції однієї накладної можуть містити однакові товари 
(при цьому ціна може бути і різною, і однаковою).

\bigskip\textit{Для деяких варіантів може знадобитись}

Недоотриману вигоду розраховувати як суму виразів
{\center (найбільша ціна відвантаження товару – ціна відвантаження товару)*кількість\par}

\bigskip
{\noindent} округлених до 2 знаків після крапки.


\par
\newpage
{\center Варіант  \No \textbf { 5 }\par}
\bigskip
%type = booking
%variant = 118
Обробляється інформація щодо відвантажених накладних.

\underline{Записи основного файлу містять поля}: вартість, номер товарної позиції у межах накладної, ціна, номер накладної, одиниця виміру, кількість, назва товару.

\underline{У допоміжному файлі наявні ключі}: найбільша ціна товарної позиції, сума вартостей.

\underline{Знайти} 3 найменш популярні товари (ті, що сумарно відвантажувалися на найменшу суму). Вивести по кожному з них інформацію:\par
{\noindent}--- на першому рядку:\par
найбільша відвантажена кількість в одній накладній (з 1 знаком після крапки),  назва товару, одиниця виміру;

{\noindent}--- на наступних рядках, починаючи з табуляції, вивести для товару всі накладні, де він відвантажувався в найбільшій кількості (по одному на рядок):\par
вартість, ціна, кількість, номер накладної

{\noindent}у такому сортуванні: ціна (за спаданням), номер накладної.

%Товари сортувати: загальна відвантажена кількість (за спаданням), назва товару, одиниця виміру.

\bigskip\textit{Обмеження}

Кількість, ціна та вартість є значеннями дійсного типу.
Решта полів є непорожніми рядками.
Рядки починатися та закінчуватися білими символами не можуть.

Номер накладної складається з 10 символів (цифри, латинські літери).
Назва товару може мати від 2 до 27 символів включно.
Може містити тільки цифри, літери, пробіли, апостроф, дефіс, підкреслення.

Кількість, ціна мають бути додатними.
Кількість вказується з рівно 1 знаками після крапки.
Ціна вказується з рівно 2 знаками після крапки.
Вартість має бути узгоджена з кількістю та ціною та записується рівно з двома десятковими знаками.
(Застосовується округлення.)

Товари однозначно ідентифікуються назвами та одиницями виміру (за наявності).

Одиниця виміру (за наявності у варіанті) може приймати значення: шт., кг, м , кор.

Кожна накладна має свій власний унікальний номер.
У накладній можуть відвантажуватися кілька товарів (не менше 1).
Товарна позиція --- рядок накладної: товар, кількість, ціна, вартість та, за наявності, одиниця виміру.
Товарні позиції кожної накладної нумеруються послідовними натуральними числами, починаючи з 1.
В одній накладній товари можуть повторюватися, 
тобто товарні позиції однієї накладної можуть містити однакові товари 
(при цьому ціна може бути і різною, і однаковою).

\bigskip\textit{Для деяких варіантів може знадобитись}

Недоотриману вигоду розраховувати як суму виразів
{\center (найбільша ціна відвантаження товару – ціна відвантаження товару)*кількість\par}

\bigskip
{\noindent} округлених до 2 знаків після крапки.


\par
\newpage
{\center Варіант  \No \textbf { 6 }\par}
\bigskip
%type = booking
%variant = 130
Обробляється інформація щодо відвантажених накладних.

\underline{Записи основного файлу містять поля}: назва товару, вартість, кількість, ціна, номер товарної позиції у межах накладної, одиниця виміру, номер накладної.

\underline{У допоміжному файлі наявні ключі}: найбільша ціна товарної позиції, кількість накладних.

\underline{Знайти} 10 накладних, в яких було відвантажено найменшу кількість різних товарів. Вивести по кожному з них інформацію:\par
{\noindent}--- на першому рядку:\par
номер накладної, кількість різних товарів у накладній, найбільша вартість товарної позиції накладної;

{\noindent}--- на наступних рядках, починаючи з табуляції, вивести для накладної всі її товарні позиції (по одному на рядок):\par
номер товарної позиції, назва товару, ціна, кількість, вартість

{\noindent}у такому сортуванні: назва товару, ціна (за спаданням).

%Накладні сортувати: кількість різних товарів, номер накладної.

\bigskip\textit{Обмеження}

Кількість, ціна та вартість є значеннями дійсного типу.
Решта полів є непорожніми рядками.
Рядки починатися та закінчуватися білими символами не можуть.

Номер накладної складається з 13 символів (цифри, латинські літери).
Назва товару може мати від 3 до 26 символів включно.
Може містити тільки цифри, літери, пробіли, апостроф, дефіс, підкреслення.

Кількість, ціна мають бути додатними.
Кількість вказується з рівно 3 знаками після крапки.
Ціна вказується з рівно 3 знаками після крапки.
Вартість має бути узгоджена з кількістю та ціною та записується рівно з двома десятковими знаками.
(Застосовується округлення.)

Товари однозначно ідентифікуються назвами та одиницями виміру (за наявності).

Одиниця виміру (за наявності у варіанті) може приймати значення: шт., кг, м , кор.

Кожна накладна має свій власний унікальний номер.
У накладній можуть відвантажуватися кілька товарів (не менше 1).
Товарна позиція --- рядок накладної: товар, кількість, ціна, вартість та, за наявності, одиниця виміру.
Товарні позиції кожної накладної нумеруються послідовними натуральними числами, починаючи з 1.
В одній накладній товари можуть повторюватися, 
тобто товарні позиції однієї накладної можуть містити однакові товари 
(при цьому ціна може бути і різною, і однаковою).

\bigskip\textit{Для деяких варіантів може знадобитись}

Недоотриману вигоду розраховувати як суму виразів
{\center (найбільша ціна відвантаження товару – ціна відвантаження товару)*кількість\par}

\bigskip
{\noindent} округлених до 2 знаків після крапки.


\par
\newpage
{\center Варіант  \No \textbf { 7 }\par}
\bigskip
%type = booking
%variant = 123
Обробляється інформація щодо відвантажених накладних.

\underline{Записи основного файлу містять поля}: назва товару, кількість, ціна, номер товарної позиції у межах накладної, вартість, номер накладної.

\underline{У допоміжному файлі наявні ключі}: середня вартість, найбільша ціна товарної позиції.

\underline{Знайти} 7 найдешевших товарів (з найменшою ціною відвантаження). Вивести по кожному з них інформацію:\par
{\noindent}--- на першому рядку:\par
назва товару, загальна відвантажена кількість, середня ціна відвантаження (з 3 знаками після крапки), найменша ціна відвантаження;

{\noindent}--- на наступних рядках, починаючи з табуляції, вивести для товару всі накладні, де він відвантажувався за ціною, що менша середньої (по одному на рядок):\par
номер накладної, кількість, ціна, вартість

{\noindent}у такому сортуванні: ціна (за спаданням), кількість, номер накладної.

%Товари сортувати: середня ціна відвантаження, назва товару.

\bigskip\textit{Обмеження}

Кількість, ціна та вартість є значеннями дійсного типу.
Решта полів є непорожніми рядками.
Рядки починатися та закінчуватися білими символами не можуть.

Номер накладної складається з 14 символів (цифри, латинські літери).
Назва товару може мати від 1 до 23 символів включно.
Може містити тільки цифри, літери, пробіли, апостроф, дефіс, підкреслення.

Кількість, ціна мають бути додатними.
Кількість вказується з рівно 2 знаками після крапки.
Ціна вказується з рівно 4 знаками після крапки.
Вартість має бути узгоджена з кількістю та ціною та записується рівно з двома десятковими знаками.
(Застосовується округлення.)

Товари однозначно ідентифікуються назвами та одиницями виміру (за наявності).

Одиниця виміру (за наявності у варіанті) може приймати значення: шт., кг, м , кор.

Кожна накладна має свій власний унікальний номер.
У накладній можуть відвантажуватися кілька товарів (не менше 1).
Товарна позиція --- рядок накладної: товар, кількість, ціна, вартість та, за наявності, одиниця виміру.
Товарні позиції кожної накладної нумеруються послідовними натуральними числами, починаючи з 1.
В одній накладній товари можуть повторюватися, 
тобто товарні позиції однієї накладної можуть містити однакові товари 
(при цьому ціна може бути і різною, і однаковою).

\bigskip\textit{Для деяких варіантів може знадобитись}

Недоотриману вигоду розраховувати як суму виразів
{\center (найбільша ціна відвантаження товару – ціна відвантаження товару)*кількість\par}

\bigskip
{\noindent} округлених до 2 знаків після крапки.


\par
\newpage
{\center Варіант  \No \textbf { 8 }\par}
\bigskip
%type = booking
%variant = 119
Обробляється інформація щодо відвантажених накладних.

\underline{Записи основного файлу містять поля}: ціна, кількість, номер товарної позиції у межах накладної, назва товару, номер накладної, вартість.

\underline{У допоміжному файлі наявні ключі}: кількість різних товарів, кількість накладних.

\underline{Знайти} 4 найбільш популярні товари (ті, що сумарно відвантажувалися на найбільшу суму). Вивести по кожному з них інформацію:\par
{\noindent}--- на першому рядку:\par
назва товару, загальна сума відвантаження (з 2 знаками після крапки), середня ціна відвантаження (з 2 знаками після крапки);

{\noindent}--- на наступних рядках, починаючи з табуляції, вивести для товару всі накладні, де він відвантажувався в найбільшій кількості (по одному на рядок):\par
номер накладної, вартість, ціна, кількість

{\noindent}у такому сортуванні: номер накладної, ціна (за спаданням).

%Товари сортувати: назва товару.

\bigskip\textit{Обмеження}

Кількість, ціна та вартість є значеннями дійсного типу.
Решта полів є непорожніми рядками.
Рядки починатися та закінчуватися білими символами не можуть.

Номер накладної складається з 13 символів (цифри, латинські літери).
Назва товару може мати від 4 до 21 символів включно.
Може містити тільки цифри, літери, пробіли, апостроф, дефіс, підкреслення.

Кількість, ціна мають бути додатними.
Кількість вказується з рівно 4 знаками після крапки.
Ціна вказується з рівно 3 знаками після крапки.
Вартість має бути узгоджена з кількістю та ціною та записується рівно з двома десятковими знаками.
(Застосовується округлення.)

Товари однозначно ідентифікуються назвами та одиницями виміру (за наявності).

Одиниця виміру (за наявності у варіанті) може приймати значення: шт., кг, м , кор.

Кожна накладна має свій власний унікальний номер.
У накладній можуть відвантажуватися кілька товарів (не менше 1).
Товарна позиція --- рядок накладної: товар, кількість, ціна, вартість та, за наявності, одиниця виміру.
Товарні позиції кожної накладної нумеруються послідовними натуральними числами, починаючи з 1.
В одній накладній товари можуть повторюватися, 
тобто товарні позиції однієї накладної можуть містити однакові товари 
(при цьому ціна може бути і різною, і однаковою).

\bigskip\textit{Для деяких варіантів може знадобитись}

Недоотриману вигоду розраховувати як суму виразів
{\center (найбільша ціна відвантаження товару – ціна відвантаження товару)*кількість\par}

\bigskip
{\noindent} округлених до 2 знаків після крапки.


\par
\newpage
{\center Варіант  \No \textbf { 9 }\par}
\bigskip
%type = booking
%variant = 121
Обробляється інформація щодо відвантажених накладних.

\underline{Записи основного файлу містять поля}: вартість, номер накладної, одиниця виміру, назва товару, кількість, номер товарної позиції у межах накладної, ціна.

\underline{У допоміжному файлі наявні ключі}: сума цін, найменша ціна товарної позиції.

\underline{Знайти} 5 найбільш популярних товарів (ті, що відвантажувалися в найбільшій кількості накладних). Вивести по кожному з них інформацію:\par
{\noindent}--- на першому рядку:\par
кількість накладних, в яких відвантажувалися товари, назва товару, одиниця виміру, сума відвантаження (з 2 знаками після крапки), найменша відпускна ціна товару (з 4 знаками після крапки);

{\noindent}--- на наступних рядках, починаючи з табуляції, вивести для товару всі накладні, де він відвантажувався менше ніж в кількості 2 одиниць (по одному на рядок):\par
номер накладної, кількість, ціна

{\noindent}у такому сортуванні: кількість (за спаданням), номер накладної.

%Товари сортувати: кількість накладних з товаром (за спаданням), назва товару, одиниця виміру.

\bigskip\textit{Обмеження}

Кількість, ціна та вартість є значеннями дійсного типу.
Решта полів є непорожніми рядками.
Рядки починатися та закінчуватися білими символами не можуть.

Номер накладної складається з 12 символів (цифри, латинські літери).
Назва товару може мати від 4 до 20 символів включно.
Може містити тільки цифри, літери, пробіли, апостроф, дефіс, підкреслення.

Кількість, ціна мають бути додатними.
Кількість вказується з рівно 3 знаками після крапки.
Ціна вказується з рівно 2 знаками після крапки.
Вартість має бути узгоджена з кількістю та ціною та записується рівно з двома десятковими знаками.
(Застосовується округлення.)

Товари однозначно ідентифікуються назвами та одиницями виміру (за наявності).

Одиниця виміру (за наявності у варіанті) може приймати значення: шт., кг, м , кор.

Кожна накладна має свій власний унікальний номер.
У накладній можуть відвантажуватися кілька товарів (не менше 1).
Товарна позиція --- рядок накладної: товар, кількість, ціна, вартість та, за наявності, одиниця виміру.
Товарні позиції кожної накладної нумеруються послідовними натуральними числами, починаючи з 1.
В одній накладній товари можуть повторюватися, 
тобто товарні позиції однієї накладної можуть містити однакові товари 
(при цьому ціна може бути і різною, і однаковою).

\bigskip\textit{Для деяких варіантів може знадобитись}

Недоотриману вигоду розраховувати як суму виразів
{\center (найбільша ціна відвантаження товару – ціна відвантаження товару)*кількість\par}

\bigskip
{\noindent} округлених до 2 знаків після крапки.


\par
\newpage
{\center Варіант  \No \textbf { 10 }\par}
\bigskip
%type = booking
%variant = 129
Обробляється інформація щодо відвантажених накладних.

\underline{Записи основного файлу містять поля}: вартість, одиниця виміру, ціна, назва товару, номер товарної позиції у межах накладної, кількість, номер накладної.

\underline{У допоміжному файлі наявні ключі}: сума кількостей, середня вартість.

\underline{Знайти} всі накладні, в яких було відвантажено більше 10 різних товарів. Вивести по кожному з них інформацію:\par
{\noindent}--- на першому рядку:\par
номер накладної, кількість різних товарів, кількість товарних позицій, середня ціна товарної позиції (з 3 знаками після крапки);

{\noindent}--- на наступних рядках, починаючи з табуляції, вивести для накладної всі товарні позиції, що відвантажувалися в ній на найменшу вартість (по одному на рядок):\par
номер товарної позиції, ціна, кількість, вартість

{\noindent}у такому сортуванні: номер товарної позиції.

%Накладні сортувати: номер накладної (за спаданням).

\bigskip\textit{Обмеження}

Кількість, ціна та вартість є значеннями дійсного типу.
Решта полів є непорожніми рядками.
Рядки починатися та закінчуватися білими символами не можуть.

Номер накладної складається з 11 символів (цифри, латинські літери).
Назва товару може мати від 3 до 22 символів включно.
Може містити тільки цифри, літери, пробіли, апостроф, дефіс, підкреслення.

Кількість, ціна мають бути додатними.
Кількість вказується з рівно 1 знаками після крапки.
Ціна вказується з рівно 4 знаками після крапки.
Вартість має бути узгоджена з кількістю та ціною та записується рівно з двома десятковими знаками.
(Застосовується округлення.)

Товари однозначно ідентифікуються назвами та одиницями виміру (за наявності).

Одиниця виміру (за наявності у варіанті) може приймати значення: шт., кг, м , кор.

Кожна накладна має свій власний унікальний номер.
У накладній можуть відвантажуватися кілька товарів (не менше 1).
Товарна позиція --- рядок накладної: товар, кількість, ціна, вартість та, за наявності, одиниця виміру.
Товарні позиції кожної накладної нумеруються послідовними натуральними числами, починаючи з 1.
В одній накладній товари можуть повторюватися, 
тобто товарні позиції однієї накладної можуть містити однакові товари 
(при цьому ціна може бути і різною, і однаковою).

\bigskip\textit{Для деяких варіантів може знадобитись}

Недоотриману вигоду розраховувати як суму виразів
{\center (найбільша ціна відвантаження товару – ціна відвантаження товару)*кількість\par}

\bigskip
{\noindent} округлених до 2 знаків після крапки.


\par
\newpage
{\center Варіант  \No \textbf { 11 }\par}
\bigskip
%type = booking
%variant = 128
Обробляється інформація щодо відвантажених накладних.

\underline{Записи основного файлу містять поля}: номер товарної позиції у межах накладної, номер накладної, назва товару, кількість, вартість, ціна.

\underline{У допоміжному файлі наявні ключі}: сума вартостей, сума цін.

\underline{Знайти} 5 накладних, в яких було відвантажено товарів на найменшу суму. Вивести по кожному з них інформацію:\par
{\noindent}--- на першому рядку:\par
середня ціна товарної позиції (з 2 знаками після крапки), загальна сума, номер накладної, кількість товарних позицій;

{\noindent}--- на наступних рядках, починаючи з табуляції, вивести для накладної всі її товарні позиції з найбільшою вартістю (по одному на рядок):\par
назва товару, ціна, кількість, вартість

{\noindent}у такому сортуванні: вартість, назва товару.

%Накладні сортувати: загальна сума, номер накладної.

\bigskip\textit{Обмеження}

Кількість, ціна та вартість є значеннями дійсного типу.
Решта полів є непорожніми рядками.
Рядки починатися та закінчуватися білими символами не можуть.

Номер накладної складається з 10 символів (цифри, латинські літери).
Назва товару може мати від 2 до 30 символів включно.
Може містити тільки цифри, літери, пробіли, апостроф, дефіс, підкреслення.

Кількість, ціна мають бути додатними.
Кількість вказується з рівно 2 знаками після крапки.
Ціна вказується з рівно 3 знаками після крапки.
Вартість має бути узгоджена з кількістю та ціною та записується рівно з двома десятковими знаками.
(Застосовується округлення.)

Товари однозначно ідентифікуються назвами та одиницями виміру (за наявності).

Одиниця виміру (за наявності у варіанті) може приймати значення: шт., кг, м , кор.

Кожна накладна має свій власний унікальний номер.
У накладній можуть відвантажуватися кілька товарів (не менше 1).
Товарна позиція --- рядок накладної: товар, кількість, ціна, вартість та, за наявності, одиниця виміру.
Товарні позиції кожної накладної нумеруються послідовними натуральними числами, починаючи з 1.
В одній накладній товари можуть повторюватися, 
тобто товарні позиції однієї накладної можуть містити однакові товари 
(при цьому ціна може бути і різною, і однаковою).

\bigskip\textit{Для деяких варіантів може знадобитись}

Недоотриману вигоду розраховувати як суму виразів
{\center (найбільша ціна відвантаження товару – ціна відвантаження товару)*кількість\par}

\bigskip
{\noindent} округлених до 2 знаків після крапки.


\par
\newpage
{\center Варіант  \No \textbf { 12 }\par}
\bigskip
%type = booking
%variant = 126
Обробляється інформація щодо відвантажених накладних.

\underline{Записи основного файлу містять поля}: номер товарної позиції у межах накладної, назва товару, вартість, одиниця виміру, ціна, кількість, номер накладної.

\underline{У допоміжному файлі наявні ключі}: найменша ціна товарної позиції, кількість записів.

\underline{Знайти} 5 накладних, в яких було відвантажено найбільшу кількість різних товарів. Вивести по кожному з них інформацію:\par
{\noindent}--- на першому рядку:\par
номер накладної, кількість товару, кількість товарних позицій, кількість різних товарів, загальна сума, середня ціна товарної позиції (з 2 знаками після крапки);

{\noindent}--- на наступних рядках, починаючи з табуляції, вивести для накладної всі товарні позиції, що відвантажувалися в ній в найбільшій кількості (по одному на рядок):\par
номер товарної позиції, ціна, кількість, вартість, назва товару, одиниця виміру

{\noindent}у такому сортуванні: номер товарної позиції.

%Накладні сортувати: номер накладної (за спаданням).

\bigskip\textit{Обмеження}

Кількість, ціна та вартість є значеннями дійсного типу.
Решта полів є непорожніми рядками.
Рядки починатися та закінчуватися білими символами не можуть.

Номер накладної складається з 15 символів (цифри, латинські літери).
Назва товару може мати від 1 до 29 символів включно.
Може містити тільки цифри, літери, пробіли, апостроф, дефіс, підкреслення.

Кількість, ціна мають бути додатними.
Кількість вказується з рівно 4 знаками після крапки.
Ціна вказується з рівно 2 знаками після крапки.
Вартість має бути узгоджена з кількістю та ціною та записується рівно з двома десятковими знаками.
(Застосовується округлення.)

Товари однозначно ідентифікуються назвами та одиницями виміру (за наявності).

Одиниця виміру (за наявності у варіанті) може приймати значення: шт., кг, м , кор.

Кожна накладна має свій власний унікальний номер.
У накладній можуть відвантажуватися кілька товарів (не менше 1).
Товарна позиція --- рядок накладної: товар, кількість, ціна, вартість та, за наявності, одиниця виміру.
Товарні позиції кожної накладної нумеруються послідовними натуральними числами, починаючи з 1.
В одній накладній товари можуть повторюватися, 
тобто товарні позиції однієї накладної можуть містити однакові товари 
(при цьому ціна може бути і різною, і однаковою).

\bigskip\textit{Для деяких варіантів може знадобитись}

Недоотриману вигоду розраховувати як суму виразів
{\center (найбільша ціна відвантаження товару – ціна відвантаження товару)*кількість\par}

\bigskip
{\noindent} округлених до 2 знаків після крапки.


\par
\newpage
{\center Варіант  \No \textbf { 13 }\par}
\bigskip
%type = booking
%variant = 127
Обробляється інформація щодо відвантажених накладних.

\underline{Записи основного файлу містять поля}: кількість, номер накладної, вартість, назва товару, ціна, номер товарної позиції у межах накладної.

\underline{У допоміжному файлі наявні ключі}: середня кількість товарних позицій у накладній, найбільша кількість товарних позицій у накладній.

\underline{Знайти} всі накладні, в яких було відвантажено рівно один товар. Вивести по кожному з них інформацію:\par
{\noindent}--- на першому рядку:\par
номер накладної, кількість товару, кількість товарних позицій, загальна сума;

{\noindent}--- на наступних рядках, починаючи з табуляції, вивести для накладної всі її товарні позиції (по одному на рядок):\par
номер товарної позиції, назва товару, ціна, кількість, вартість

{\noindent}у такому сортуванні: назва товару, ціна.

%Накладні сортувати: кількість товару (за спаданням), номер накладної.

\bigskip\textit{Обмеження}

Кількість, ціна та вартість є значеннями дійсного типу.
Решта полів є непорожніми рядками.
Рядки починатися та закінчуватися білими символами не можуть.

Номер накладної складається з 15 символів (цифри, латинські літери).
Назва товару може мати від 1 до 26 символів включно.
Може містити тільки цифри, літери, пробіли, апостроф, дефіс, підкреслення.

Кількість, ціна мають бути додатними.
Кількість вказується з рівно 2 знаками після крапки.
Ціна вказується з рівно 3 знаками після крапки.
Вартість має бути узгоджена з кількістю та ціною та записується рівно з двома десятковими знаками.
(Застосовується округлення.)

Товари однозначно ідентифікуються назвами та одиницями виміру (за наявності).

Одиниця виміру (за наявності у варіанті) може приймати значення: шт., кг, м , кор.

Кожна накладна має свій власний унікальний номер.
У накладній можуть відвантажуватися кілька товарів (не менше 1).
Товарна позиція --- рядок накладної: товар, кількість, ціна, вартість та, за наявності, одиниця виміру.
Товарні позиції кожної накладної нумеруються послідовними натуральними числами, починаючи з 1.
В одній накладній товари можуть повторюватися, 
тобто товарні позиції однієї накладної можуть містити однакові товари 
(при цьому ціна може бути і різною, і однаковою).

\bigskip\textit{Для деяких варіантів може знадобитись}

Недоотриману вигоду розраховувати як суму виразів
{\center (найбільша ціна відвантаження товару – ціна відвантаження товару)*кількість\par}

\bigskip
{\noindent} округлених до 2 знаків після крапки.


\par
\newpage
{\center Варіант  \No \textbf { 14 }\par}
\bigskip
%type =     booking
%variant = 116
Обробляється інформація щодо відвантажених накладних.

\underline{Записи основного файлу містять поля}: номер накладної, ціна, вартість, назва товару, кількість, номер товарної позиції у межах накладної.

\underline{У допоміжному файлі наявні ключі}: середня вартість, кількість записів.

\underline{Знайти} 4 найменш популярні товари (ті, що сумарно відвантажувалися в найменшій кількості). Вивести по кожному з них інформацію:\par
{\noindent}--- на першому рядку:\par
назва товару, загальна відвантажена кількість, найбільша вартість відвантаження (з 2 знаками після крапки);

{\noindent}--- на наступних рядках, починаючи з табуляції, вивести для товару всі накладні, де його було відвантажено більше 2 одиниць (по одному на рядок):\par
номер накладної, кількість, ціна, вартість

{\noindent}у такому сортуванні: номер накладної, кількість (за спаданням).

%Накладні сортувати: загальна відвантажена кількість (за спаданням), назва.

\bigskip\textit{Обмеження}

Кількість, ціна та вартість є значеннями дійсного типу.
Решта полів є непорожніми рядками.
Рядки починатися та закінчуватися білими символами не можуть.

Номер накладної складається з 13 символів (цифри, латинські літери).
Назва товару може мати від 4 до 20 символів включно.
Може містити тільки цифри, літери, пробіли, апостроф, дефіс, підкреслення.

Кількість, ціна мають бути додатними.
Кількість вказується з рівно 4 знаками після крапки.
Ціна вказується з рівно 4 знаками після крапки.
Вартість має бути узгоджена з кількістю та ціною та записується рівно з двома десятковими знаками.
(Застосовується округлення.)

Товари однозначно ідентифікуються назвами та одиницями виміру (за наявності).

Одиниця виміру (за наявності у варіанті) може приймати значення: шт., кг, м , кор.

Кожна накладна має свій власний унікальний номер.
У накладній можуть відвантажуватися кілька товарів (не менше 1).
Товарна позиція --- рядок накладної: товар, кількість, ціна, вартість та, за наявності, одиниця виміру.
Товарні позиції кожної накладної нумеруються послідовними натуральними числами, починаючи з 1.
В одній накладній товари можуть повторюватися, 
тобто товарні позиції однієї накладної можуть містити однакові товари 
(при цьому ціна може бути і різною, і однаковою).

\bigskip\textit{Для деяких варіантів може знадобитись}

Недоотриману вигоду розраховувати як суму виразів
{\center (найбільша ціна відвантаження товару – ціна відвантаження товару)*кількість\par}

\bigskip
{\noindent} округлених до 2 знаків після крапки.


\par
\newpage
{\center Варіант  \No \textbf { 15 }\par}
\bigskip
%type = booking
%variant = 131
Обробляється інформація щодо відвантажених накладних.

\underline{Записи основного файлу містять поля}: кількість, номер накладної, назва товару, ціна, номер товарної позиції у межах накладної, вартість.

\underline{У допоміжному файлі наявні ключі}: сума цін, сума кількостей.

\underline{Знайти} 5 накладних, в яких було відвантажено товарів на найбільшу суму. Вивести по кожному з них інформацію:\par
{\noindent}--- на першому рядку:\par
номер накладної, кількість товару, кількість товарних позицій, загальна сума, середня вартість товарів у накладній (з 2 знаками після крапки);

{\noindent}--- на наступних рядках, починаючи з табуляції, вивести для накладної всі її товарні позиції, де відвантажувалося менше 1 одиниці товару (по одному на рядок):\par
назва товару, ціна, кількість, вартість, номер товарної позиції

{\noindent}у такому сортуванні: назва товару, кількість, ціна.

%Накладні сортувати: середня вартість товарів у накладній (за спаданням), номер накладної.

\bigskip\textit{Обмеження}

Кількість, ціна та вартість є значеннями дійсного типу.
Решта полів є непорожніми рядками.
Рядки починатися та закінчуватися білими символами не можуть.

Номер накладної складається з 11 символів (цифри, латинські літери).
Назва товару може мати від 2 до 30 символів включно.
Може містити тільки цифри, літери, пробіли, апостроф, дефіс, підкреслення.

Кількість, ціна мають бути додатними.
Кількість вказується з рівно 3 знаками після крапки.
Ціна вказується з рівно 2 знаками після крапки.
Вартість має бути узгоджена з кількістю та ціною та записується рівно з двома десятковими знаками.
(Застосовується округлення.)

Товари однозначно ідентифікуються назвами та одиницями виміру (за наявності).

Одиниця виміру (за наявності у варіанті) може приймати значення: шт., кг, м , кор.

Кожна накладна має свій власний унікальний номер.
У накладній можуть відвантажуватися кілька товарів (не менше 1).
Товарна позиція --- рядок накладної: товар, кількість, ціна, вартість та, за наявності, одиниця виміру.
Товарні позиції кожної накладної нумеруються послідовними натуральними числами, починаючи з 1.
В одній накладній товари можуть повторюватися, 
тобто товарні позиції однієї накладної можуть містити однакові товари 
(при цьому ціна може бути і різною, і однаковою).

\bigskip\textit{Для деяких варіантів може знадобитись}

Недоотриману вигоду розраховувати як суму виразів
{\center (найбільша ціна відвантаження товару – ціна відвантаження товару)*кількість\par}

\bigskip
{\noindent} округлених до 2 знаків після крапки.


\par
\newpage
{\center Варіант  \No \textbf { 16 }\par}
\bigskip
%type = booking
%variant = 120
Обробляється інформація щодо відвантажених накладних.

\underline{Записи основного файлу містять поля}: ціна, назва товару, вартість, кількість, номер товарної позиції у межах накладної, номер накладної.

\underline{У допоміжному файлі наявні ключі}: сума вартостей, кількість різних товарів.

\underline{Знайти} три найбільш популярні товари (ті, що сумарно відвантажувалися в найбільшій кількості). Вивести по кожному з них інформацію:\par
{\noindent}--- на першому рядку:\par
назва товару, сума відвантаження (з 2 знаками після крапки), найбільша відвантажена кількість в одній накладній (з 1 знаком після крапки);

{\noindent}--- на наступних рядках, починаючи з табуляції, вивести для товару всі накладні, де він відвантажувався більш ніж на 1000 (по одному на рядок):\par
вартість, ціна, кількість, номер накладної

{\noindent}у такому сортуванні: кількість, вартість, ціна, номер накладної.

%Товари сортувати: назва товару (за спаданням).

\bigskip\textit{Обмеження}

Кількість, ціна та вартість є значеннями дійсного типу.
Решта полів є непорожніми рядками.
Рядки починатися та закінчуватися білими символами не можуть.

Номер накладної складається з 14 символів (цифри, латинські літери).
Назва товару може мати від 3 до 21 символів включно.
Може містити тільки цифри, літери, пробіли, апостроф, дефіс, підкреслення.

Кількість, ціна мають бути додатними.
Кількість вказується з рівно 1 знаками після крапки.
Ціна вказується з рівно 4 знаками після крапки.
Вартість має бути узгоджена з кількістю та ціною та записується рівно з двома десятковими знаками.
(Застосовується округлення.)

Товари однозначно ідентифікуються назвами та одиницями виміру (за наявності).

Одиниця виміру (за наявності у варіанті) може приймати значення: шт., кг, м , кор.

Кожна накладна має свій власний унікальний номер.
У накладній можуть відвантажуватися кілька товарів (не менше 1).
Товарна позиція --- рядок накладної: товар, кількість, ціна, вартість та, за наявності, одиниця виміру.
Товарні позиції кожної накладної нумеруються послідовними натуральними числами, починаючи з 1.
В одній накладній товари можуть повторюватися, 
тобто товарні позиції однієї накладної можуть містити однакові товари 
(при цьому ціна може бути і різною, і однаковою).

\bigskip\textit{Для деяких варіантів може знадобитись}

Недоотриману вигоду розраховувати як суму виразів
{\center (найбільша ціна відвантаження товару – ціна відвантаження товару)*кількість\par}

\bigskip
{\noindent} округлених до 2 знаків після крапки.


\par
\newpage
{\center Варіант  \No \textbf { 17 }\par}
\bigskip
%type = booking
%variant = 131
Обробляється інформація щодо відвантажених накладних.

\underline{Записи основного файлу містять поля}: кількість, номер накладної, назва товару, ціна, номер товарної позиції у межах накладної, вартість.

\underline{У допоміжному файлі наявні ключі}: сума цін, сума кількостей.

\underline{Знайти} 5 накладних, в яких було відвантажено товарів на найбільшу суму. Вивести по кожному з них інформацію:\par
{\noindent}--- на першому рядку:\par
номер накладної, кількість товару, кількість товарних позицій, загальна сума, середня вартість товарів у накладній (з 2 знаками після крапки);

{\noindent}--- на наступних рядках, починаючи з табуляції, вивести для накладної всі її товарні позиції, де відвантажувалося менше 1 одиниці товару (по одному на рядок):\par
назва товару, ціна, кількість, вартість, номер товарної позиції

{\noindent}у такому сортуванні: назва товару, кількість, ціна.

%Накладні сортувати: середня вартість товарів у накладній (за спаданням), номер накладної.

\bigskip\textit{Обмеження}

Кількість, ціна та вартість є значеннями дійсного типу.
Решта полів є непорожніми рядками.
Рядки починатися та закінчуватися білими символами не можуть.

Номер накладної складається з 10 символів (цифри, латинські літери).
Назва товару може мати від 1 до 22 символів включно.
Може містити тільки цифри, літери, пробіли, апостроф, дефіс, підкреслення.

Кількість, ціна мають бути додатними.
Кількість вказується з рівно 4 знаками після крапки.
Ціна вказується з рівно 3 знаками після крапки.
Вартість має бути узгоджена з кількістю та ціною та записується рівно з двома десятковими знаками.
(Застосовується округлення.)

Товари однозначно ідентифікуються назвами та одиницями виміру (за наявності).

Одиниця виміру (за наявності у варіанті) може приймати значення: шт., кг, м , кор.

Кожна накладна має свій власний унікальний номер.
У накладній можуть відвантажуватися кілька товарів (не менше 1).
Товарна позиція --- рядок накладної: товар, кількість, ціна, вартість та, за наявності, одиниця виміру.
Товарні позиції кожної накладної нумеруються послідовними натуральними числами, починаючи з 1.
В одній накладній товари можуть повторюватися, 
тобто товарні позиції однієї накладної можуть містити однакові товари 
(при цьому ціна може бути і різною, і однаковою).

\bigskip\textit{Для деяких варіантів може знадобитись}

Недоотриману вигоду розраховувати як суму виразів
{\center (найбільша ціна відвантаження товару – ціна відвантаження товару)*кількість\par}

\bigskip
{\noindent} округлених до 2 знаків після крапки.


\par
\newpage
{\center Варіант  \No \textbf { 18 }\par}
\bigskip
%type = booking
%variant = 124
Обробляється інформація щодо відвантажених накладних.

\underline{Записи основного файлу містять поля}: номер накладної, кількість, назва товару, вартість, ціна, номер товарної позиції у межах накладної.

\underline{У допоміжному файлі наявні ключі}: найбільша ціна товарної позиції, сума вартостей.

\underline{Знайти} 5 найдорожчих товарів (з найбільшою середньою ціною відвантаження). Вивести по кожному з них інформацію:\par
{\noindent}--- на першому рядку:\par
назва товару, загальна відвантажена кількість, середня ціна відвантаження (з 2 знаками після крапки), найменша ціна відвантаження,  найбільша ціна відвантаження;

{\noindent}--- на наступних рядках, починаючи з табуляції, вивести для товару всі накладні, де він відвантажувався за ненайбільшою своєю ціною (по одному на рядок):\par
ціна, кількість, ціна, вартість, номер накладної

{\noindent}у такому сортуванні: вартість (за спаданням), ціна(за спаданням), номер накладної.

%Товари сортувати: середня ціна відвантаження (за спаданням), назва товару.

\bigskip\textit{Обмеження}

Кількість, ціна та вартість є значеннями дійсного типу.
Решта полів є непорожніми рядками.
Рядки починатися та закінчуватися білими символами не можуть.

Номер накладної складається з 12 символів (цифри, латинські літери).
Назва товару може мати від 3 до 24 символів включно.
Може містити тільки цифри, літери, пробіли, апостроф, дефіс, підкреслення.

Кількість, ціна мають бути додатними.
Кількість вказується з рівно 3 знаками після крапки.
Ціна вказується з рівно 2 знаками після крапки.
Вартість має бути узгоджена з кількістю та ціною та записується рівно з двома десятковими знаками.
(Застосовується округлення.)

Товари однозначно ідентифікуються назвами та одиницями виміру (за наявності).

Одиниця виміру (за наявності у варіанті) може приймати значення: шт., кг, м , кор.

Кожна накладна має свій власний унікальний номер.
У накладній можуть відвантажуватися кілька товарів (не менше 1).
Товарна позиція --- рядок накладної: товар, кількість, ціна, вартість та, за наявності, одиниця виміру.
Товарні позиції кожної накладної нумеруються послідовними натуральними числами, починаючи з 1.
В одній накладній товари можуть повторюватися, 
тобто товарні позиції однієї накладної можуть містити однакові товари 
(при цьому ціна може бути і різною, і однаковою).

\bigskip\textit{Для деяких варіантів може знадобитись}

Недоотриману вигоду розраховувати як суму виразів
{\center (найбільша ціна відвантаження товару – ціна відвантаження товару)*кількість\par}

\bigskip
{\noindent} округлених до 2 знаків після крапки.


\par
\newpage
{\center Варіант  \No \textbf { 19 }\par}
\bigskip
%type = booking
%variant = 128
Обробляється інформація щодо відвантажених накладних.

\underline{Записи основного файлу містять поля}: номер товарної позиції у межах накладної, номер накладної, назва товару, кількість, вартість, ціна.

\underline{У допоміжному файлі наявні ключі}: сума вартостей, сума цін.

\underline{Знайти} 5 накладних, в яких було відвантажено товарів на найменшу суму. Вивести по кожному з них інформацію:\par
{\noindent}--- на першому рядку:\par
середня ціна товарної позиції (з 2 знаками після крапки), загальна сума, номер накладної, кількість товарних позицій;

{\noindent}--- на наступних рядках, починаючи з табуляції, вивести для накладної всі її товарні позиції з найбільшою вартістю (по одному на рядок):\par
назва товару, ціна, кількість, вартість

{\noindent}у такому сортуванні: вартість, назва товару.

%Накладні сортувати: загальна сума, номер накладної.

\bigskip\textit{Обмеження}

Кількість, ціна та вартість є значеннями дійсного типу.
Решта полів є непорожніми рядками.
Рядки починатися та закінчуватися білими символами не можуть.

Номер накладної складається з 11 символів (цифри, латинські літери).
Назва товару може мати від 4 до 28 символів включно.
Може містити тільки цифри, літери, пробіли, апостроф, дефіс, підкреслення.

Кількість, ціна мають бути додатними.
Кількість вказується з рівно 2 знаками після крапки.
Ціна вказується з рівно 4 знаками після крапки.
Вартість має бути узгоджена з кількістю та ціною та записується рівно з двома десятковими знаками.
(Застосовується округлення.)

Товари однозначно ідентифікуються назвами та одиницями виміру (за наявності).

Одиниця виміру (за наявності у варіанті) може приймати значення: шт., кг, м , кор.

Кожна накладна має свій власний унікальний номер.
У накладній можуть відвантажуватися кілька товарів (не менше 1).
Товарна позиція --- рядок накладної: товар, кількість, ціна, вартість та, за наявності, одиниця виміру.
Товарні позиції кожної накладної нумеруються послідовними натуральними числами, починаючи з 1.
В одній накладній товари можуть повторюватися, 
тобто товарні позиції однієї накладної можуть містити однакові товари 
(при цьому ціна може бути і різною, і однаковою).

\bigskip\textit{Для деяких варіантів може знадобитись}

Недоотриману вигоду розраховувати як суму виразів
{\center (найбільша ціна відвантаження товару – ціна відвантаження товару)*кількість\par}

\bigskip
{\noindent} округлених до 2 знаків після крапки.


\par
\newpage
{\center Варіант  \No \textbf { 20 }\par}
\bigskip
%type = booking
%variant = 120
Обробляється інформація щодо відвантажених накладних.

\underline{Записи основного файлу містять поля}: ціна, назва товару, вартість, кількість, номер товарної позиції у межах накладної, номер накладної.

\underline{У допоміжному файлі наявні ключі}: сума вартостей, кількість різних товарів.

\underline{Знайти} три найбільш популярні товари (ті, що сумарно відвантажувалися в найбільшій кількості). Вивести по кожному з них інформацію:\par
{\noindent}--- на першому рядку:\par
назва товару, сума відвантаження (з 2 знаками після крапки), найбільша відвантажена кількість в одній накладній (з 1 знаком після крапки);

{\noindent}--- на наступних рядках, починаючи з табуляції, вивести для товару всі накладні, де він відвантажувався більш ніж на 1000 (по одному на рядок):\par
вартість, ціна, кількість, номер накладної

{\noindent}у такому сортуванні: кількість, вартість, ціна, номер накладної.

%Товари сортувати: назва товару (за спаданням).

\bigskip\textit{Обмеження}

Кількість, ціна та вартість є значеннями дійсного типу.
Решта полів є непорожніми рядками.
Рядки починатися та закінчуватися білими символами не можуть.

Номер накладної складається з 13 символів (цифри, латинські літери).
Назва товару може мати від 2 до 23 символів включно.
Може містити тільки цифри, літери, пробіли, апостроф, дефіс, підкреслення.

Кількість, ціна мають бути додатними.
Кількість вказується з рівно 1 знаками після крапки.
Ціна вказується з рівно 2 знаками після крапки.
Вартість має бути узгоджена з кількістю та ціною та записується рівно з двома десятковими знаками.
(Застосовується округлення.)

Товари однозначно ідентифікуються назвами та одиницями виміру (за наявності).

Одиниця виміру (за наявності у варіанті) може приймати значення: шт., кг, м , кор.

Кожна накладна має свій власний унікальний номер.
У накладній можуть відвантажуватися кілька товарів (не менше 1).
Товарна позиція --- рядок накладної: товар, кількість, ціна, вартість та, за наявності, одиниця виміру.
Товарні позиції кожної накладної нумеруються послідовними натуральними числами, починаючи з 1.
В одній накладній товари можуть повторюватися, 
тобто товарні позиції однієї накладної можуть містити однакові товари 
(при цьому ціна може бути і різною, і однаковою).

\bigskip\textit{Для деяких варіантів може знадобитись}

Недоотриману вигоду розраховувати як суму виразів
{\center (найбільша ціна відвантаження товару – ціна відвантаження товару)*кількість\par}

\bigskip
{\noindent} округлених до 2 знаків після крапки.


\par
\newpage
{\center Варіант  \No \textbf { 21 }\par}
\bigskip
%type =     booking
%variant = 116
Обробляється інформація щодо відвантажених накладних.

\underline{Записи основного файлу містять поля}: номер накладної, ціна, вартість, назва товару, кількість, номер товарної позиції у межах накладної.

\underline{У допоміжному файлі наявні ключі}: середня вартість, кількість записів.

\underline{Знайти} 4 найменш популярні товари (ті, що сумарно відвантажувалися в найменшій кількості). Вивести по кожному з них інформацію:\par
{\noindent}--- на першому рядку:\par
назва товару, загальна відвантажена кількість, найбільша вартість відвантаження (з 2 знаками після крапки);

{\noindent}--- на наступних рядках, починаючи з табуляції, вивести для товару всі накладні, де його було відвантажено більше 2 одиниць (по одному на рядок):\par
номер накладної, кількість, ціна, вартість

{\noindent}у такому сортуванні: номер накладної, кількість (за спаданням).

%Накладні сортувати: загальна відвантажена кількість (за спаданням), назва.

\bigskip\textit{Обмеження}

Кількість, ціна та вартість є значеннями дійсного типу.
Решта полів є непорожніми рядками.
Рядки починатися та закінчуватися білими символами не можуть.

Номер накладної складається з 15 символів (цифри, латинські літери).
Назва товару може мати від 4 до 25 символів включно.
Може містити тільки цифри, літери, пробіли, апостроф, дефіс, підкреслення.

Кількість, ціна мають бути додатними.
Кількість вказується з рівно 2 знаками після крапки.
Ціна вказується з рівно 3 знаками після крапки.
Вартість має бути узгоджена з кількістю та ціною та записується рівно з двома десятковими знаками.
(Застосовується округлення.)

Товари однозначно ідентифікуються назвами та одиницями виміру (за наявності).

Одиниця виміру (за наявності у варіанті) може приймати значення: шт., кг, м , кор.

Кожна накладна має свій власний унікальний номер.
У накладній можуть відвантажуватися кілька товарів (не менше 1).
Товарна позиція --- рядок накладної: товар, кількість, ціна, вартість та, за наявності, одиниця виміру.
Товарні позиції кожної накладної нумеруються послідовними натуральними числами, починаючи з 1.
В одній накладній товари можуть повторюватися, 
тобто товарні позиції однієї накладної можуть містити однакові товари 
(при цьому ціна може бути і різною, і однаковою).

\bigskip\textit{Для деяких варіантів може знадобитись}

Недоотриману вигоду розраховувати як суму виразів
{\center (найбільша ціна відвантаження товару – ціна відвантаження товару)*кількість\par}

\bigskip
{\noindent} округлених до 2 знаків після крапки.


\par
\newpage
{\center Варіант  \No \textbf { 22 }\par}
\bigskip
%type = booking
%variant = 122
Обробляється інформація щодо відвантажених накладних.

\underline{Записи основного файлу містять поля}: назва товару, кількість, ціна, номер товарної позиції у межах накладної, номер накладної, вартість, одиниця виміру.

\underline{У допоміжному файлі наявні ключі}: кількість накладних, кількість записів.

\underline{Знайти} 5 найдорожчих товарів (з найбільшою ціною відвантаження). Вивести по кожному з них інформацію:\par
{\noindent}--- на першому рядку:\par
найбільша ціна відвантаження, найменша ціна відвантаження, назва товару, одиниця виміру, загальна відвантажена кількість;

{\noindent}--- на наступних рядках, починаючи з табуляції, вивести для товару всі накладні, де він відвантажувався за ціною, що відрізняється від найбільшої (по одному на рядок):\par
номер накладної, кількість, ціна, недоотримана вигода (з 1 знаком після крапки)

{\noindent}у такому сортуванні: номер накладної.

%Товари сортувати: різниця максимальної та мінімальної відпускних цін за спаданням, назва товару, одиниця виміру.

\bigskip\textit{Обмеження}

Кількість, ціна та вартість є значеннями дійсного типу.
Решта полів є непорожніми рядками.
Рядки починатися та закінчуватися білими символами не можуть.

Номер накладної складається з 14 символів (цифри, латинські літери).
Назва товару може мати від 3 до 27 символів включно.
Може містити тільки цифри, літери, пробіли, апостроф, дефіс, підкреслення.

Кількість, ціна мають бути додатними.
Кількість вказується з рівно 1 знаками після крапки.
Ціна вказується з рівно 4 знаками після крапки.
Вартість має бути узгоджена з кількістю та ціною та записується рівно з двома десятковими знаками.
(Застосовується округлення.)

Товари однозначно ідентифікуються назвами та одиницями виміру (за наявності).

Одиниця виміру (за наявності у варіанті) може приймати значення: шт., кг, м , кор.

Кожна накладна має свій власний унікальний номер.
У накладній можуть відвантажуватися кілька товарів (не менше 1).
Товарна позиція --- рядок накладної: товар, кількість, ціна, вартість та, за наявності, одиниця виміру.
Товарні позиції кожної накладної нумеруються послідовними натуральними числами, починаючи з 1.
В одній накладній товари можуть повторюватися, 
тобто товарні позиції однієї накладної можуть містити однакові товари 
(при цьому ціна може бути і різною, і однаковою).

\bigskip\textit{Для деяких варіантів може знадобитись}

Недоотриману вигоду розраховувати як суму виразів
{\center (найбільша ціна відвантаження товару – ціна відвантаження товару)*кількість\par}

\bigskip
{\noindent} округлених до 2 знаків після крапки.


\par
\newpage
{\center Варіант  \No \textbf { 23 }\par}
\bigskip
%type = booking
%variant = 117
Обробляється інформація щодо відвантажених накладних.

\underline{Записи основного файлу містять поля}: ціна, номер накладної, назву товару, вартість, одиниця виміру, номер товарної позиції у межах накладної, кількість.

\underline{У допоміжному файлі наявні ключі}: сума цін, найменша ціна товарної позиції.

\underline{Знайти} 6 найменш популярних товарів (ті, що відвантажувалися в найменшій кількості накладних). Вивести по кожному з них інформацію:\par
{\noindent}--- на першому рядку:\par
назва товару, одиниця виміру, загальна відвантажена кількість, найменша ціна відвантаження (з 2 знаками після крапки);

{\noindent}--- на наступних рядках, починаючи з табуляції, вивести для товару всі накладні, де він відвантажувався за найбільшою ціною (по одному на рядок):\par
ціна, кількість, номер накладної

{\noindent}у такому сортуванні: ціна, номер накладної.

%Товари сортувати: загальна відвантажена кількість, найменша ціна відвантаження, назва товару, одиниця виміру.

\bigskip\textit{Обмеження}

Кількість, ціна та вартість є значеннями дійсного типу.
Решта полів є непорожніми рядками.
Рядки починатися та закінчуватися білими символами не можуть.

Номер накладної складається з 12 символів (цифри, латинські літери).
Назва товару може мати від 2 до 26 символів включно.
Може містити тільки цифри, літери, пробіли, апостроф, дефіс, підкреслення.

Кількість, ціна мають бути додатними.
Кількість вказується з рівно 3 знаками після крапки.
Ціна вказується з рівно 3 знаками після крапки.
Вартість має бути узгоджена з кількістю та ціною та записується рівно з двома десятковими знаками.
(Застосовується округлення.)

Товари однозначно ідентифікуються назвами та одиницями виміру (за наявності).

Одиниця виміру (за наявності у варіанті) може приймати значення: шт., кг, м , кор.

Кожна накладна має свій власний унікальний номер.
У накладній можуть відвантажуватися кілька товарів (не менше 1).
Товарна позиція --- рядок накладної: товар, кількість, ціна, вартість та, за наявності, одиниця виміру.
Товарні позиції кожної накладної нумеруються послідовними натуральними числами, починаючи з 1.
В одній накладній товари можуть повторюватися, 
тобто товарні позиції однієї накладної можуть містити однакові товари 
(при цьому ціна може бути і різною, і однаковою).

\bigskip\textit{Для деяких варіантів може знадобитись}

Недоотриману вигоду розраховувати як суму виразів
{\center (найбільша ціна відвантаження товару – ціна відвантаження товару)*кількість\par}

\bigskip
{\noindent} округлених до 2 знаків після крапки.


\par
\newpage
{\center Варіант  \No \textbf { 24 }\par}
\bigskip
%type = booking
%variant = 123
Обробляється інформація щодо відвантажених накладних.

\underline{Записи основного файлу містять поля}: назва товару, кількість, ціна, номер товарної позиції у межах накладної, вартість, номер накладної.

\underline{У допоміжному файлі наявні ключі}: середня вартість, найбільша ціна товарної позиції.

\underline{Знайти} 7 найдешевших товарів (з найменшою ціною відвантаження). Вивести по кожному з них інформацію:\par
{\noindent}--- на першому рядку:\par
назва товару, загальна відвантажена кількість, середня ціна відвантаження (з 3 знаками після крапки), найменша ціна відвантаження;

{\noindent}--- на наступних рядках, починаючи з табуляції, вивести для товару всі накладні, де він відвантажувався за ціною, що менша середньої (по одному на рядок):\par
номер накладної, кількість, ціна, вартість

{\noindent}у такому сортуванні: ціна (за спаданням), кількість, номер накладної.

%Товари сортувати: середня ціна відвантаження, назва товару.

\bigskip\textit{Обмеження}

Кількість, ціна та вартість є значеннями дійсного типу.
Решта полів є непорожніми рядками.
Рядки починатися та закінчуватися білими символами не можуть.

Номер накладної складається з 10 символів (цифри, латинські літери).
Назва товару може мати від 1 до 24 символів включно.
Може містити тільки цифри, літери, пробіли, апостроф, дефіс, підкреслення.

Кількість, ціна мають бути додатними.
Кількість вказується з рівно 4 знаками після крапки.
Ціна вказується з рівно 2 знаками після крапки.
Вартість має бути узгоджена з кількістю та ціною та записується рівно з двома десятковими знаками.
(Застосовується округлення.)

Товари однозначно ідентифікуються назвами та одиницями виміру (за наявності).

Одиниця виміру (за наявності у варіанті) може приймати значення: шт., кг, м , кор.

Кожна накладна має свій власний унікальний номер.
У накладній можуть відвантажуватися кілька товарів (не менше 1).
Товарна позиція --- рядок накладної: товар, кількість, ціна, вартість та, за наявності, одиниця виміру.
Товарні позиції кожної накладної нумеруються послідовними натуральними числами, починаючи з 1.
В одній накладній товари можуть повторюватися, 
тобто товарні позиції однієї накладної можуть містити однакові товари 
(при цьому ціна може бути і різною, і однаковою).

\bigskip\textit{Для деяких варіантів може знадобитись}

Недоотриману вигоду розраховувати як суму виразів
{\center (найбільша ціна відвантаження товару – ціна відвантаження товару)*кількість\par}

\bigskip
{\noindent} округлених до 2 знаків після крапки.


\par
\newpage
{\center Варіант  \No \textbf { 25 }\par}
\bigskip
%type = booking
%variant = 125
Обробляється інформація щодо відвантажених накладних.

\underline{Записи основного файлу містять поля}: кількість, одиниця виміру, ціна, назва товару, номер товарної позиції у межах накладної, вартість, номер накладної.

\underline{У допоміжному файлі наявні ключі}: кількість записів, кількість різних товарів.

\underline{Знайти} 7 найдешевших товарів (з найменшою середньою ціною відвантаження). Вивести по кожному з них інформацію:\par
{\noindent}--- на першому рядку:\par
назва товару, одиниця виміру, кількість накладних, де відвантажувався товар, кількість накладних, де він відвантажувався за своєю максимальною ціною, середня ціна відвантаження (з 2 знаками після крапки);

{\noindent}--- на наступних рядках, починаючи з табуляції, вивести для товару всі накладні, де він відвантажувався за своєю мінімальною ціною (по одному на рядок):\par
номер накладної, ціна, кількість, ціна, вартість

{\noindent}у такому сортуванні: номер накладної, ціна.

%Товари сортувати: кількість накладних, де відвантажувався товар, назва товару, одиниця виміру.

\bigskip\textit{Обмеження}

Кількість, ціна та вартість є значеннями дійсного типу.
Решта полів є непорожніми рядками.
Рядки починатися та закінчуватися білими символами не можуть.

Номер накладної складається з 12 символів (цифри, латинські літери).
Назва товару може мати від 1 до 29 символів включно.
Може містити тільки цифри, літери, пробіли, апостроф, дефіс, підкреслення.

Кількість, ціна мають бути додатними.
Кількість вказується з рівно 3 знаками після крапки.
Ціна вказується з рівно 4 знаками після крапки.
Вартість має бути узгоджена з кількістю та ціною та записується рівно з двома десятковими знаками.
(Застосовується округлення.)

Товари однозначно ідентифікуються назвами та одиницями виміру (за наявності).

Одиниця виміру (за наявності у варіанті) може приймати значення: шт., кг, м , кор.

Кожна накладна має свій власний унікальний номер.
У накладній можуть відвантажуватися кілька товарів (не менше 1).
Товарна позиція --- рядок накладної: товар, кількість, ціна, вартість та, за наявності, одиниця виміру.
Товарні позиції кожної накладної нумеруються послідовними натуральними числами, починаючи з 1.
В одній накладній товари можуть повторюватися, 
тобто товарні позиції однієї накладної можуть містити однакові товари 
(при цьому ціна може бути і різною, і однаковою).

\bigskip\textit{Для деяких варіантів може знадобитись}

Недоотриману вигоду розраховувати як суму виразів
{\center (найбільша ціна відвантаження товару – ціна відвантаження товару)*кількість\par}

\bigskip
{\noindent} округлених до 2 знаків після крапки.


\par
\newpage
{\center Варіант  \No \textbf { 26 }\par}
\bigskip
%type = booking
%variant = 119
Обробляється інформація щодо відвантажених накладних.

\underline{Записи основного файлу містять поля}: ціна, кількість, номер товарної позиції у межах накладної, назва товару, номер накладної, вартість.

\underline{У допоміжному файлі наявні ключі}: кількість різних товарів, кількість накладних.

\underline{Знайти} 4 найбільш популярні товари (ті, що сумарно відвантажувалися на найбільшу суму). Вивести по кожному з них інформацію:\par
{\noindent}--- на першому рядку:\par
назва товару, загальна сума відвантаження (з 2 знаками після крапки), середня ціна відвантаження (з 2 знаками після крапки);

{\noindent}--- на наступних рядках, починаючи з табуляції, вивести для товару всі накладні, де він відвантажувався в найбільшій кількості (по одному на рядок):\par
номер накладної, вартість, ціна, кількість

{\noindent}у такому сортуванні: номер накладної, ціна (за спаданням).

%Товари сортувати: назва товару.

\bigskip\textit{Обмеження}

Кількість, ціна та вартість є значеннями дійсного типу.
Решта полів є непорожніми рядками.
Рядки починатися та закінчуватися білими символами не можуть.

Номер накладної складається з 14 символів (цифри, латинські літери).
Назва товару може мати від 4 до 22 символів включно.
Може містити тільки цифри, літери, пробіли, апостроф, дефіс, підкреслення.

Кількість, ціна мають бути додатними.
Кількість вказується з рівно 4 знаками після крапки.
Ціна вказується з рівно 2 знаками після крапки.
Вартість має бути узгоджена з кількістю та ціною та записується рівно з двома десятковими знаками.
(Застосовується округлення.)

Товари однозначно ідентифікуються назвами та одиницями виміру (за наявності).

Одиниця виміру (за наявності у варіанті) може приймати значення: шт., кг, м , кор.

Кожна накладна має свій власний унікальний номер.
У накладній можуть відвантажуватися кілька товарів (не менше 1).
Товарна позиція --- рядок накладної: товар, кількість, ціна, вартість та, за наявності, одиниця виміру.
Товарні позиції кожної накладної нумеруються послідовними натуральними числами, починаючи з 1.
В одній накладній товари можуть повторюватися, 
тобто товарні позиції однієї накладної можуть містити однакові товари 
(при цьому ціна може бути і різною, і однаковою).

\bigskip\textit{Для деяких варіантів може знадобитись}

Недоотриману вигоду розраховувати як суму виразів
{\center (найбільша ціна відвантаження товару – ціна відвантаження товару)*кількість\par}

\bigskip
{\noindent} округлених до 2 знаків після крапки.


\par
\newpage
{\center Варіант  \No \textbf { 27 }\par}
\bigskip
%type = booking
%variant = 121
Обробляється інформація щодо відвантажених накладних.

\underline{Записи основного файлу містять поля}: вартість, номер накладної, одиниця виміру, назва товару, кількість, номер товарної позиції у межах накладної, ціна.

\underline{У допоміжному файлі наявні ключі}: сума цін, найменша ціна товарної позиції.

\underline{Знайти} 5 найбільш популярних товарів (ті, що відвантажувалися в найбільшій кількості накладних). Вивести по кожному з них інформацію:\par
{\noindent}--- на першому рядку:\par
кількість накладних, в яких відвантажувалися товари, назва товару, одиниця виміру, сума відвантаження (з 2 знаками після крапки), найменша відпускна ціна товару (з 4 знаками після крапки);

{\noindent}--- на наступних рядках, починаючи з табуляції, вивести для товару всі накладні, де він відвантажувався менше ніж в кількості 2 одиниць (по одному на рядок):\par
номер накладної, кількість, ціна

{\noindent}у такому сортуванні: кількість (за спаданням), номер накладної.

%Товари сортувати: кількість накладних з товаром (за спаданням), назва товару, одиниця виміру.

\bigskip\textit{Обмеження}

Кількість, ціна та вартість є значеннями дійсного типу.
Решта полів є непорожніми рядками.
Рядки починатися та закінчуватися білими символами не можуть.

Номер накладної складається з 11 символів (цифри, латинські літери).
Назва товару може мати від 3 до 23 символів включно.
Може містити тільки цифри, літери, пробіли, апостроф, дефіс, підкреслення.

Кількість, ціна мають бути додатними.
Кількість вказується з рівно 1 знаками після крапки.
Ціна вказується з рівно 3 знаками після крапки.
Вартість має бути узгоджена з кількістю та ціною та записується рівно з двома десятковими знаками.
(Застосовується округлення.)

Товари однозначно ідентифікуються назвами та одиницями виміру (за наявності).

Одиниця виміру (за наявності у варіанті) може приймати значення: шт., кг, м , кор.

Кожна накладна має свій власний унікальний номер.
У накладній можуть відвантажуватися кілька товарів (не менше 1).
Товарна позиція --- рядок накладної: товар, кількість, ціна, вартість та, за наявності, одиниця виміру.
Товарні позиції кожної накладної нумеруються послідовними натуральними числами, починаючи з 1.
В одній накладній товари можуть повторюватися, 
тобто товарні позиції однієї накладної можуть містити однакові товари 
(при цьому ціна може бути і різною, і однаковою).

\bigskip\textit{Для деяких варіантів може знадобитись}

Недоотриману вигоду розраховувати як суму виразів
{\center (найбільша ціна відвантаження товару – ціна відвантаження товару)*кількість\par}

\bigskip
{\noindent} округлених до 2 знаків після крапки.


\par
\newpage
{\center Варіант  \No \textbf { 28 }\par}
\bigskip
%type = booking
%variant = 127
Обробляється інформація щодо відвантажених накладних.

\underline{Записи основного файлу містять поля}: кількість, номер накладної, вартість, назва товару, ціна, номер товарної позиції у межах накладної.

\underline{У допоміжному файлі наявні ключі}: середня кількість товарних позицій у накладній, найбільша кількість товарних позицій у накладній.

\underline{Знайти} всі накладні, в яких було відвантажено рівно один товар. Вивести по кожному з них інформацію:\par
{\noindent}--- на першому рядку:\par
номер накладної, кількість товару, кількість товарних позицій, загальна сума;

{\noindent}--- на наступних рядках, починаючи з табуляції, вивести для накладної всі її товарні позиції (по одному на рядок):\par
номер товарної позиції, назва товару, ціна, кількість, вартість

{\noindent}у такому сортуванні: назва товару, ціна.

%Накладні сортувати: кількість товару (за спаданням), номер накладної.

\bigskip\textit{Обмеження}

Кількість, ціна та вартість є значеннями дійсного типу.
Решта полів є непорожніми рядками.
Рядки починатися та закінчуватися білими символами не можуть.

Номер накладної складається з 15 символів (цифри, латинські літери).
Назва товару може мати від 2 до 25 символів включно.
Може містити тільки цифри, літери, пробіли, апостроф, дефіс, підкреслення.

Кількість, ціна мають бути додатними.
Кількість вказується з рівно 2 знаками після крапки.
Ціна вказується з рівно 4 знаками після крапки.
Вартість має бути узгоджена з кількістю та ціною та записується рівно з двома десятковими знаками.
(Застосовується округлення.)

Товари однозначно ідентифікуються назвами та одиницями виміру (за наявності).

Одиниця виміру (за наявності у варіанті) може приймати значення: шт., кг, м , кор.

Кожна накладна має свій власний унікальний номер.
У накладній можуть відвантажуватися кілька товарів (не менше 1).
Товарна позиція --- рядок накладної: товар, кількість, ціна, вартість та, за наявності, одиниця виміру.
Товарні позиції кожної накладної нумеруються послідовними натуральними числами, починаючи з 1.
В одній накладній товари можуть повторюватися, 
тобто товарні позиції однієї накладної можуть містити однакові товари 
(при цьому ціна може бути і різною, і однаковою).

\bigskip\textit{Для деяких варіантів може знадобитись}

Недоотриману вигоду розраховувати як суму виразів
{\center (найбільша ціна відвантаження товару – ціна відвантаження товару)*кількість\par}

\bigskip
{\noindent} округлених до 2 знаків після крапки.


\par
\newpage
{\center Варіант  \No \textbf { 29 }\par}
\bigskip
%type = booking
%variant = 126
Обробляється інформація щодо відвантажених накладних.

\underline{Записи основного файлу містять поля}: номер товарної позиції у межах накладної, назва товару, вартість, одиниця виміру, ціна, кількість, номер накладної.

\underline{У допоміжному файлі наявні ключі}: найменша ціна товарної позиції, кількість записів.

\underline{Знайти} 5 накладних, в яких було відвантажено найбільшу кількість різних товарів. Вивести по кожному з них інформацію:\par
{\noindent}--- на першому рядку:\par
номер накладної, кількість товару, кількість товарних позицій, кількість різних товарів, загальна сума, середня ціна товарної позиції (з 2 знаками після крапки);

{\noindent}--- на наступних рядках, починаючи з табуляції, вивести для накладної всі товарні позиції, що відвантажувалися в ній в найбільшій кількості (по одному на рядок):\par
номер товарної позиції, ціна, кількість, вартість, назва товару, одиниця виміру

{\noindent}у такому сортуванні: номер товарної позиції.

%Накладні сортувати: номер накладної (за спаданням).

\bigskip\textit{Обмеження}

Кількість, ціна та вартість є значеннями дійсного типу.
Решта полів є непорожніми рядками.
Рядки починатися та закінчуватися білими символами не можуть.

Номер накладної складається з 10 символів (цифри, латинські літери).
Назва товару може мати від 2 до 30 символів включно.
Може містити тільки цифри, літери, пробіли, апостроф, дефіс, підкреслення.

Кількість, ціна мають бути додатними.
Кількість вказується з рівно 4 знаками після крапки.
Ціна вказується з рівно 2 знаками після крапки.
Вартість має бути узгоджена з кількістю та ціною та записується рівно з двома десятковими знаками.
(Застосовується округлення.)

Товари однозначно ідентифікуються назвами та одиницями виміру (за наявності).

Одиниця виміру (за наявності у варіанті) може приймати значення: шт., кг, м , кор.

Кожна накладна має свій власний унікальний номер.
У накладній можуть відвантажуватися кілька товарів (не менше 1).
Товарна позиція --- рядок накладної: товар, кількість, ціна, вартість та, за наявності, одиниця виміру.
Товарні позиції кожної накладної нумеруються послідовними натуральними числами, починаючи з 1.
В одній накладній товари можуть повторюватися, 
тобто товарні позиції однієї накладної можуть містити однакові товари 
(при цьому ціна може бути і різною, і однаковою).

\bigskip\textit{Для деяких варіантів може знадобитись}

Недоотриману вигоду розраховувати як суму виразів
{\center (найбільша ціна відвантаження товару – ціна відвантаження товару)*кількість\par}

\bigskip
{\noindent} округлених до 2 знаків після крапки.


\par
\newpage
{\center Варіант  \No \textbf { 30 }\par}
\bigskip
%type = booking
%variant = 130
Обробляється інформація щодо відвантажених накладних.

\underline{Записи основного файлу містять поля}: назва товару, вартість, кількість, ціна, номер товарної позиції у межах накладної, одиниця виміру, номер накладної.

\underline{У допоміжному файлі наявні ключі}: найбільша ціна товарної позиції, кількість накладних.

\underline{Знайти} 10 накладних, в яких було відвантажено найменшу кількість різних товарів. Вивести по кожному з них інформацію:\par
{\noindent}--- на першому рядку:\par
номер накладної, кількість різних товарів у накладній, найбільша вартість товарної позиції накладної;

{\noindent}--- на наступних рядках, починаючи з табуляції, вивести для накладної всі її товарні позиції (по одному на рядок):\par
номер товарної позиції, назва товару, ціна, кількість, вартість

{\noindent}у такому сортуванні: назва товару, ціна (за спаданням).

%Накладні сортувати: кількість різних товарів, номер накладної.

\bigskip\textit{Обмеження}

Кількість, ціна та вартість є значеннями дійсного типу.
Решта полів є непорожніми рядками.
Рядки починатися та закінчуватися білими символами не можуть.

Номер накладної складається з 13 символів (цифри, латинські літери).
Назва товару може мати від 3 до 28 символів включно.
Може містити тільки цифри, літери, пробіли, апостроф, дефіс, підкреслення.

Кількість, ціна мають бути додатними.
Кількість вказується з рівно 2 знаками після крапки.
Ціна вказується з рівно 3 знаками після крапки.
Вартість має бути узгоджена з кількістю та ціною та записується рівно з двома десятковими знаками.
(Застосовується округлення.)

Товари однозначно ідентифікуються назвами та одиницями виміру (за наявності).

Одиниця виміру (за наявності у варіанті) може приймати значення: шт., кг, м , кор.

Кожна накладна має свій власний унікальний номер.
У накладній можуть відвантажуватися кілька товарів (не менше 1).
Товарна позиція --- рядок накладної: товар, кількість, ціна, вартість та, за наявності, одиниця виміру.
Товарні позиції кожної накладної нумеруються послідовними натуральними числами, починаючи з 1.
В одній накладній товари можуть повторюватися, 
тобто товарні позиції однієї накладної можуть містити однакові товари 
(при цьому ціна може бути і різною, і однаковою).

\bigskip\textit{Для деяких варіантів може знадобитись}

Недоотриману вигоду розраховувати як суму виразів
{\center (найбільша ціна відвантаження товару – ціна відвантаження товару)*кількість\par}

\bigskip
{\noindent} округлених до 2 знаків після крапки.


\par
\newpage
\end{document}
